% 第6章: 予備校XとYの合格率 (上) と受験者数 (下). 説明のために作った数字.
%   X: 難関大 100人中30人 (30%), それ以外 400人中320人 (80%), 全体 500人中350人 (70%)
%   Y: 難関大 400人中140人 (35%), それ以外 100人中85人 (85%), 全体 500人中225人 (45%)
\begin{tikzpicture}[font=\footnotesize]
  \begin{axis}[
    name=rate,
    width=0.62\linewidth, height=4.4cm, scale only axis=false,
    ybar, bar width=12pt,
    ymin=0, ymax=100, ytick={0,20,40,60,80,100},
    ylabel={合格率 (\%)},
    xmin=-0.6, xmax=2.6,
    xtick={0,1,2},
    xticklabels={全体,難関大,それ以外の大学},
    x tick label style={font=\footnotesize},
    axis x line*=bottom, axis y line*=left,
    legend style={draw=none, at={(0.5,1.04)}, anchor=south, font=\footnotesize, legend columns=2, /tikz/every even column/.append style={column sep=1em}},
    legend image code/.code={\draw[#1] (0cm,-0.1cm) rectangle (0.3cm,0.1cm);},
    enlarge x limits=false,
    nodes near coords, nodes near coords style={font=\scriptsize},
  ]
    \addplot[draw=black, fill=black!15] coordinates {(0,70) (1,30) (2,80)};
    \addplot[draw=black, fill=black!65] coordinates {(0,45) (1,35) (2,85)};
    \legend{予備校X, 予備校Y}
    \draw[black!40] (axis cs:0.5,0) -- (axis cs:0.5,100);
  \end{axis}
  \begin{axis}[
    at={(rate.south)}, anchor=north, yshift=-1.1cm,
    width=0.62\linewidth, height=3.0cm, scale only axis=false,
    ybar, bar width=12pt,
    ymin=0, ymax=500, ytick={0,200,400},
    ylabel={受験者数 (人)},
    xmin=-0.6, xmax=2.6,
    xtick={1,2},
    xticklabels={難関大,それ以外の大学},
    x tick label style={font=\footnotesize},
    axis x line*=bottom, axis y line*=left,
    enlarge x limits=false,
  ]
    \addplot[draw=black, fill=black!15] coordinates {(1,100) (2,400)};
    \addplot[draw=black, fill=black!65] coordinates {(1,400) (2,100)};
  \end{axis}
\end{tikzpicture}
